\documentclass[11pt]{article}

\usepackage[T1]{fontenc}
\usepackage{amsmath}
\usepackage{amssymb}
\usepackage{amsthm}
\usepackage{mathtools}
\usepackage{microtype}
\usepackage{hyperref}

\hypersetup{
    hidelinks,
    pdftitle={The Integrality Gap in Robin's Inequality with a Fixed Number of Prime Factors},
    pdfauthor={Logan Bell}
}
\setlength{\emergencystretch}{2em}
\numberwithin{equation}{section}

\newcommand{\reals}{\mathbf{R}}
\newcommand{\integers}{\mathbf{Z}}
\newcommand{\eps}{\varepsilon}
\DeclareMathOperator*{\argmax}{argmax}
\newcommand{\Igap}{\Delta}
\newcommand{\rstar}{r^\star}

\theoremstyle{plain}
\newtheorem{theorem}{Theorem}[section]
\newtheorem{proposition}[theorem]{Proposition}
\newtheorem{lemma}[theorem]{Lemma}
\newtheorem{corollary}[theorem]{Corollary}

\theoremstyle{remark}
\newtheorem{remark}[theorem]{Remark}

\title{The Integrality Gap in Robin's Inequality\\
with a Fixed Number of Prime Factors}
\author{Logan Bell \\ \texttt{lmbell@stanford.edu}}
\date{\today}

\begin{document}

\maketitle

\begin{abstract}
Robin's inequality is equivalent to the Riemann hypothesis. We study its
largest normalized divisor sum among integers with exactly \(m\) distinct
prime factors, together with the relaxation obtained by allowing real
exponents. A maximizer uses the first \(m\) primes. One scalar equation
determines the relaxed optimizer, while sorting exponent increments by
marginal gain per unit log-size determines every integer optimizer and
gives a finite computation. We show that each integer exponent is a floor
or ceiling of its relaxed value and evaluate the sum of the resulting
rounding losses. If \(g_{\integers}(m)\) and \(g_{\reals}(m)\) are the two
optimal log values, then
\[
    g_{\reals}(m)-g_{\integers}(m)
    =
    \frac{2\sqrt2+o(1)}
    {\sqrt{p_m}\log p_m}.
\]
Here \(p_m\) is the \(m\)th prime. The two equal leading contributions
come from downward and upward rounding near the transition between
exponents one and two. The same marginal equation appears in the
Axler--Nicolas comparison of size-constrained maxima and explains its
matching coefficient. Under the Riemann hypothesis, the relaxation
eventually violates Robin's inequality while the integer optimum satisfies
it. If the hypothesis fails, two-sided Mertens excursions dominate the
gap. Thus the Riemann hypothesis is equivalent to Robin's inequality
holding on every sufficiently large fixed-\(\omega\) slice.
\end{abstract}

\section{Introduction}
\label{s-introduction}

Let \(\sigma(n)\) be the sum of the positive divisors of \(n\), and let
\(\gamma\) be Euler's constant. Gr\"onwall \cite{gronwall1913} proved
\[
    \limsup_{n\to\infty}
    \frac{\sigma(n)}{e^\gamma n\log\log n}
    =
    1.
\]
Robin \cite{robin1984} sharpened this maximal-order statement to a
pointwise criterion. The Riemann hypothesis holds if and only if
\[
    \sigma(n)<e^\gamma n\log\log n
    \qquad (n\geq5041).
\]
We maximize the normalized ratio over integers with a fixed number of
distinct prime factors. On each such slice the extremal problem becomes
a finite optimization over the exponents, and we solve it exactly in
both integer and real variables.

Write \(\omega(n)\) for the number of distinct prime factors of \(n\).
For each \(m\), let \(g_{\integers}(m)\) be the largest log-normalized
divisor sum among integers \(n\geq5041\) with \(\omega(n)=m\). A support
exchange shows that, for \(m\geq6\), every maximizer uses the first \(m\)
primes. Allowing their exponents to be real numbers at least one gives a
relaxation with value \(g_{\reals}(m)\). We study
\[
    \Igap(m)=g_{\reals}(m)-g_{\integers}(m).
\]

\paragraph{Main result.}
If \(p_m\) is the \(m\)th prime, then
\begin{equation}
\label{e-integrality-gap-intro}
    \Igap(m)
    =
    \frac{2\sqrt2+o(1)}
         {\sqrt{p_m}\log p_m}.
\end{equation}
The result is unconditional. The real optimizer is determined by one
scalar equation. Every integer optimizer is a prefix of the exponent
increments ordered by Nicolas's score \(F(p,k)\), with equal scores
taken as one block. This gives a finite
computation with exact score comparisons. Each integer exponent is one of
the two integers adjacent to its relaxed counterpart. The resulting local
rounding sum transfers to a prime integral, where only the second-exponent
range contributes at first order. The two sides of its rounding switch
each contribute \(\sqrt2/(\sqrt{p_m}\log p_m)\).

\paragraph{Robin's inequality.}
The identity
\(-g_{\integers}(m)=\Igap(m)-g_{\reals}(m)\) separates the rounding loss
from the analytic Mertens term. Under the Riemann hypothesis, Nicolas's
integrated explicit formula makes the relaxed value eventually positive
and bounds its normalized coefficient above by
\(4+\gamma-\log(4\pi)\). This is smaller than the gap coefficient
\(2\sqrt2\), so the integer optimum is eventually below Robin's bound. If
the Riemann hypothesis fails, Nicolas's two-sided Mertens excursions
dominate \eqref{e-integrality-gap-intro}. It follows that the hypothesis
is equivalent to Robin's inequality holding on every sufficiently large
fixed-\(\omega\) slice.

\paragraph{Scope.}
We solve the fixed-support relaxation by one scalar equation, characterize
every integer optimizer as a score-ordered prefix, prove one-step
proximity, and evaluate the resulting rounding loss. Combining this
unconditional calculation with Nicolas's integrated explicit formula
gives the fixed-slice criterion. Nicolas's CA score and shell estimates
also place the calculation in the classical theory. In the problem with
\(\omega(n)\leq m\), first-exponent increments become optional and the
same score ordering characterizes the optimizers. A finite verification
settles the slices with \(m\leq5\), where the cutoff \(n\geq5041\)
blocks the support exchange. The integer
\(10080=2^{5}3^{2}5\cdot7\) holds the bounded-support maximum for
\(6\leq m\leq29\), and for \(m\geq30\) the maximizers are score-ordered
prefixes.

\subsection{Relation to previous work}
\label{s-related-work}

\paragraph{Superabundant and colossally abundant numbers.}
Superabundant (SA) numbers are the image of the record map
\[
    X\mapsto\argmax_{n\leq X}\frac{\sigma(n)}n,
\]
and colossally abundant (CA) numbers are the image of the map
\[
    \eps\mapsto\argmax_{n\geq1}\frac{\sigma(n)}{n^{1+\eps}},
\]
defined for \(\eps>0\).
Ramanujan, Alaoglu and Erd\H{o}s, Nicolas, and Robin developed this
description and the associated benefit methods
\cite{ramanujan-hcn-1997,alaoglu1944,nicolas-1975-asterisque,
nicolas-1976-asterisque,robin-1980,robin-1983-rairo}.
Nicolas's marginal score \(F(p,k)\) governs whether the exponent of \(p\)
reaches \(k\) \cite[Section~4.1]{nicolas-2022-sigma}.

The score \(F(p,k)\) is also the value-to-log-size ratio in our knapsack
ordering. Fixing \(\omega(n)=m\) makes the first-exponent increments
\(F(p,1)\) mandatory for \(p\leq p_m\), while the optional increments
with \(k\geq2\) retain their classical order. The resulting optimizer is
CA and SA relative to its prescribed support, but need not be a global
CA or SA number. The bounded-support problem restores the support
decisions.

\paragraph{Classical Robin asymptotics.}
Under the Riemann hypothesis, Ramanujan's expansion for generalized
superior highly composite numbers, in the corrected version published by
Nicolas and Robin, already contains the normalized constant
\(2\sqrt2-2\) and the zeta-zero term that appear in the classical Robin
margin
\cite[formula~(382) and Note~71]{ramanujan-hcn-1997}. Nicolas gives a
modern effective treatment for CA numbers and then for all integers
\cite[Proposition~5.1 and Theorem~1.1]{nicolas-2022-sigma}. The
coefficient \(2\sqrt2-2\) and the zero sum in the final integer margin are
therefore not new here.

These results establish the classical integer margin, not the
fixed-support relaxation or its integrality gap. Our optimizer
characterization and rounding formula give that unconditional
gap. Nicolas's analytic estimates enter only afterward, when the gap is
combined with the relaxed value.

\paragraph{The Axler--Nicolas comparison.}
Axler and Nicolas \cite{axler-nicolas-2023} prove
\[
    \frac{\displaystyle\max_{n\leq X}n/\varphi(n)}
         {\displaystyle\max_{n\leq X}\sigma(n)/n}
    =
    1+
    \frac{2\sqrt2}
         {\sqrt{\log X}\log\log X}
    +
    O\left(
        \frac1{\sqrt{\log X}(\log\log X)^2}
    \right),
\]
where \(\varphi\) is Euler's totient function. Nicolas obtains a full
expansion \cite[Theorem~1.2]{nicolas-2025}.
Their size-constrained problem and our fixed-support penalized problem have
different feasible sets, so neither theorem implies the other. The
shared second-exponent score equation explains the same leading
coefficient. The second terms have different meanings in the two
problems. Section~\ref{s-score-connection} gives the precise comparison.

\paragraph{Other reductions.}
The first-prime support exchange is standard in work on possible
counterexamples to Robin's inequality
\cite[Section~5, Lemma~9]{choie2007}. Other work restricts a least
counterexample or reduces Robin's criterion to record subsequences
\cite{akbary-friggstad-2009,nazardonyavi-yakubovich-2014}. We instead
maximize separately on each level set of \(\omega(n)\). This makes the
first-exponent increments for the first \(m\) primes mandatory and leaves
the exponent lattice as the discrete part of the problem.

\paragraph{Outline.}
Section~\ref{s-fixed-slice} formulates the problem.
Section~\ref{s-optimizers} gives the scalar relaxed optimizer and the
score-ordered prefix theorem.
Section~\ref{s-integrality-gap} proves one-step proximity and evaluates
the rounding loss. Its final subsection relates the calculation
to CA shells and the Axler--Nicolas comparison.
Sections~\ref{s-relaxed-excess} and \ref{s-robin-margin} add Nicolas's
analytic input and derive the fixed-slice criterion.
Section~\ref{s-bounded-support} treats \(\omega(n)\leq m\).

\section{The fixed-\texorpdfstring{\(\omega\)}{omega} extremal problem}
\label{s-fixed-slice}

We normalize Robin's inequality by
\[
    G(n)
    =
    \frac{\sigma(n)}
    {e^\gamma n\log\log n}.
\]
Thus Robin's inequality is \(G(n)<1\) for \(n\geq5041\)
\cite{robin1984}. For \(m\geq1\), define the fixed-\(\omega\) extremal
value
\[
    g_{\integers}(m)
    =
    \sup_{\substack{n\geq5041\\ \omega(n)=m}}\log G(n).
\]
A negative value of \(g_{\integers}(m)\) means that Robin's inequality
holds throughout the slice. For \(m\geq6\), the attainment proved below
makes the converse true as well.

We write \(p_i\) for the \(i\)th prime, and use
\[
    \pi(x)=\#\{p\leq x\},
    \qquad
    \theta(x)=\sum_{p\leq x}\log p.
\]
We use \(\integers_{++}=\{1,2,\ldots\}\).
The first reduction removes the choice of prime support.

\begin{proposition}[Support reduction]
\label{prop-support}
For every \(m\geq6\), the supremum defining
\(g_{\integers}(m)\) is a maximum. Every maximizer has the form
\[
    n=\prod_{i=1}^{m}p_i^{r_i},
    \qquad
    r\in\integers_{++}^{m}.
\]
\end{proposition}

\begin{proof}
Suppose
\[
    n=\prod_{i=1}^{m}q_i^{r_i},
    \qquad
    q_1<\cdots<q_m,
    \qquad
    r_i\geq1.
\]
For a fixed exponent \(a\), the local factor
\[
    1+q^{-1}+\cdots+q^{-a}
\]
is strictly decreasing in \(q\). We replace \(q_i\) by \(p_i\) for every
\(i\). This cannot decrease \(\sigma(n)/n\) or increase \(n\). Since
\(\log\log n\) is increasing on the range in question, the replacement
cannot decrease \(G(n)\).

For \(m\geq6\), the replacement remains feasible because
\[
    \prod_{i=1}^{m}p_i
    \geq
    \prod_{i=1}^{6}p_i
    =
    30030
    >
    5041.
\]
If some \(q_i>p_i\), the corresponding local factor increases strictly,
while the denominator does not increase. Hence \(G(n)\) increases
strictly, and every maximizer uses \(p_1,\ldots,p_m\).

We now prove attainment. On this fixed support,
\(\sigma(n)/n\) is bounded above by
\[
    \prod_{i=1}^{m}\frac{p_i}{p_i-1},
\]
whereas \(\log\log n\to\infty\) when any exponent tends to infinity.
Hence \(\log G(n)\to-\infty\) whenever an exponent tends to infinity.
Only finitely many exponent vectors can
lie above the value at \(r_i=1\), so the supremum is attained.
\end{proof}

For \(m\leq5\) the product of the first \(m\) primes is smaller than
\(5041\), so the replacement in the proof can leave the range
\(n\geq5041\), and the support need not reduce to the first \(m\)
primes. A finite verification settles these slices.

\begin{proposition}[Small slices]
\label{prop-small-slices}
For \(1\leq m\leq5\), the supremum defining \(g_{\integers}(m)\) is a
maximum, attained uniquely at
\[
    2^{13},\qquad
    2^{6}3^{4},\qquad
    2^{5}3^{2}5^{2},\qquad
    2^{5}3^{2}5\cdot7,\qquad
    2^{4}3^{2}5\cdot7\cdot11,
\]
respectively. Each maximum is negative, and
\(g_{\integers}(4)=\log G(10080)=-0.0142\ldots\) exceeds the other
four.
\end{proposition}

\begin{proof}
The local factor of \(\sigma(n)/n\) at a prime \(p\) is smaller than
\(p/(p-1)\), and \(p/(p-1)\) decreases in \(p\), so \(\omega(n)=m\) gives
\(\sigma(n)/n<\prod_{i=1}^{m}p_i/(p_i-1)\). Hence
\[
    \log G(n)
    <
    \sum_{i=1}^{m}\log\frac{p_i}{p_i-1}
    -\gamma-\log\log\log n.
\]
The bound decreases in \(n\), and at \(n=10^7\) it is already below the
maximum claimed for the corresponding \(m\). A verification over the
integers in \([5041,10^7]\) with \(\omega(n)=m\) completes the proof.
\end{proof}

Robin's inequality therefore holds on every slice with \(m\leq5\). The
values are not monotone in \(m\), since
\(g_{\integers}(5)=\log G(55440)=-0.0168\ldots\) falls below the
\(m=4\) value. Section~\ref{s-bounded-support} shows that \(10080\)
holds the bounded-support maximum for \(6\leq m\leq29\).

\paragraph{Real-exponent relaxation.}
For the rest of the paper, \(m\geq6\) and the support is
\((p_1,\ldots,p_m)\). For \(r\in[1,\infty)^m\), we write
\[
    S(r)=\sum_{i=1}^{m}r_i\log p_i.
\]
For \(r\in\integers_{++}^{m}\), we have
\[
    \frac{\sigma(\prod_i p_i^{r_i})}
    {\prod_i p_i^{r_i}}
    =
    \prod_{i=1}^{m}
    \frac{p_i-p_i^{-r_i}}{p_i-1},
\]
so the real-exponent relaxation is obtained by using the expression on
the right for every \(r\in[1,\infty)^m\). Its value is
\begin{equation}
\label{e-objective}
    g_{\reals}(m)
    =
    \sup_{r\in[1,\infty)^m}
    \left(
        \sum_{i=1}^{m}\log(p_i-p_i^{-r_i})
        -\gamma
        -\sum_{i=1}^{m}\log(p_i-1)
        -\log\log S(r)
    \right).
\end{equation}
This objective agrees with \(\log G(\prod_i p_i^{r_i})\) at every
positive integer exponent vector.
Since the integer exponent vectors form a subset of its feasible set,
\[
    g_{\reals}(m)\geq g_{\integers}(m).
\]

The integrality gap is
\[
    \Igap(m)=g_{\reals}(m)-g_{\integers}(m).
\]
We refer to \(g_{\reals}(m)\) as the relaxed excess and to
\(-g_{\integers}(m)\) as the fixed-slice Robin margin. A positive relaxed
excess means that the real-exponent relaxation violates Robin's bound. A
positive Robin margin means that every integer in the slice satisfies the
bound.

Thus, for every \(m\geq6\),
\[
    -g_{\integers}(m)
    =
    \Igap(m)-g_{\reals}(m).
\]
We estimate the two terms separately. The gap is unconditional and comes
from the exponent lattice. The
relaxed excess is a normalized Mertens error and carries the dependence on
the zeros of the zeta function.

\section{The two optimizers}
\label{s-optimizers}

The relaxed and integer problems both reduce to one-dimensional
constructions. A scalar equation gives the relaxed optimizer. A
fractional-knapsack ordering gives every integer optimizer and a finite
algorithm for computing it.

\subsection{The relaxed optimizer}
\label{s-relaxed-scalar}

For \(T>0\), define the log-size of the coordinate response by
\[
    s_m(T)
    =
    \sum_{i=1}^{m}
    \max\{\log p_i,\log(1+T)-\log p_i\}.
\]

\begin{proposition}[Relaxed optimizer]
\label{prop-scalar-reduction}
The relaxed problem has a unique maximizer \(\rstar\). There is a unique
solution \(T^\star>0\) of
\[
    T=s_m(T)\log s_m(T),
\]
and
\begin{equation}
\label{e-relaxed-response}
    r_i^\star
    =
    \max\left\{
        1,\,
        \frac{\log(1+T^\star)}{\log p_i}-1
    \right\}.
\end{equation}
In particular,
\[
    S^\star:=S(\rstar)=s_m(T^\star),
    \qquad
    T^\star=S^\star\log S^\star.
\]
\end{proposition}

\begin{proof}
The local logarithmic terms in \eqref{e-objective} are bounded above,
whereas \(\log\log S(r)\) tends to infinity whenever a coordinate of
\(r\) does. A maximum therefore exists. At any maximum, put
\(T=S(r)\log S(r)\). Since
\[
    \frac{d}{du}\log(p_i-p_i^{-u})
    =
    \frac{\log p_i}{p_i^{u+1}-1},
\]
the one-sided first-order conditions give
\[
    p_i^{r_i+1}-1=T
    \quad\hbox{if \(r_i>1\)},\qquad
    p_i^2-1\geq T
    \quad\hbox{if \(r_i=1\)}.
\]
These conditions are equivalent to \eqref{e-relaxed-response}, and the
size of that response is \(s_m(T)\).

It remains to show that the scalar equation has one solution. Set
\[
    h(T)=T-s_m(T)\log s_m(T).
\]
We have \(h(0)=-\theta(p_m)\log\theta(p_m)<0\), while
\(s_m(T)\leq\theta(p_m)+m\log(1+T)\), so \(h(T)\) tends to infinity.
Every root satisfies \(T\geq\theta(p_m)\log\theta(p_m)\).

Away from the switching points, let \(k\) be the number of primes with
\(p_i^2<1+T\). If \(k=0\), then \(h'(T)=1\). Otherwise,
\[
    h'(T)
    =
    1-\frac{k(1+\log s_m(T))}{1+T}.
\]
The \(k\)th prime is at least \(k+1\), so \(1+T>(k+1)^2\). Also
\[
    s_m(T)
    =
    \theta(p_m)
    +
    \sum_{p_i^2<1+T}
    \bigl(\log(1+T)-2\log p_i\bigr)
    <2(1+T).
\]
Indeed, on the root interval \(\theta(p_m)<1+T\), while
\(k<\sqrt{1+T}\) and \(\log(1+T)<\sqrt{1+T}\).
The elementary inequality
\[
    k\bigl(1+\log(2u)\bigr)<u,
    \qquad u>(k+1)^2,
\]
now gives \(h'(T)>0\). At the left endpoint the inequality reduces to
\[
    k+1+\frac1k-\log2-2\log(k+1)>0.
\]
This is immediate for \(k=1,2\) and increasing for \(k\geq2\).
Continuity across the switching points proves uniqueness of the root.
Every global maximizer satisfies the first-order conditions, so the
response at this root is the unique maximizer.
\end{proof}

\subsection{The integer optimizer}
\label{s-integer-prefix}

For \(x>1\) and an integer \(k\geq1\), define the standard CA score
\cite[Section~4.1]{nicolas-2022-sigma}
\begin{equation}
\label{e-ca-score}
    F(x,k)
    =
    \frac1{\log x}
    \log\left(1+\frac1{x+x^2+\cdots+x^k}\right).
\end{equation}
Raising the exponent of \(p\) from \(k-1\) to \(k\) changes the log
divisor ratio by
\begin{equation}
\label{e-score-dictionary}
    \log\left(
        \frac{\sigma(p^k)/p^k}
             {\sigma(p^{k-1})/p^{k-1}}
    \right)
    =
    \log p\,F(p,k).
\end{equation}
Equivalently,
\begin{equation}
\label{e-score-integral}
    F(p,k)
    =
    \int_{k-1}^{k}\frac{du}{p^{u+1}-1}.
\end{equation}
This shows that the score decreases strictly in both arguments and that
\[
    F(p,k)\leq\frac1{p^k-1}.
\]
In particular, only finitely many eligible scores exceed any positive
level.
The fixed-support problem includes the increments \((p,k)\) with
\(p\leq p_m\) and \(k\geq2\). Their log-size weights are \(\log p\),
and their scores are ordered as above.

Group equal scores and order the resulting blocks by decreasing score.
Starting from the primorial \(\prod_{p\leq p_m}p\), let \(N_j\) be the
integer obtained after the first \(j\) blocks. We refer to these
integers as score-ordered prefixes. Every such prefix
is a valid exponent vector because each predecessor in a prime chain has
a larger score.

\begin{theorem}[Prefix characterization]
\label{thm-integer-prefix}
Every integer maximizer is one of the prefix integers \(N_j\). If \(N_m^\star\)
is any maximizer and
\[
    L_m=\log N_m^\star,
    \qquad
    \eps_m=\frac1{L_m\log L_m},
\]
then an increment \((p,k)\), \(k\geq2\), is selected exactly when
\begin{equation}
\label{e-integer-threshold}
    F(p,k)>\eps_m.
\end{equation}
No eligible score equals \(\eps_m\), and no maximizer cuts a tie block.
Every maximizing prefix also satisfies
\begin{equation}
\label{e-prefix-bound}
    \log N_m^\star\leq\theta(p_m)^{\zeta(2)}.
\end{equation}
\end{theorem}

\begin{proof}
Consider the piecewise-linear curve through the points
\[
    \left(
        \log N_j,\,
        \log\frac{\sigma(N_j)}{N_j}
    \right).
\]
Its slopes are the ordered scores. This curve is a fractional-knapsack
upper bound. Let \(\mathcal E\) be any finite selection of increments
with total extra log-size \(z\). On an interval with score \(c\), let
\(\mathcal P\) be the set of all increments in earlier blocks. Then
\[
\begin{aligned}
    \sum_{(p,k)\in\mathcal E}\log p\,F(p,k)
    &=
    cz+
    \sum_{(p,k)\in\mathcal E}
    \log p\bigl(F(p,k)-c\bigr)\\
    &\leq
    cz+
    \sum_{(p,k)\in\mathcal P}
    \log p\bigl(F(p,k)-c\bigr).
\end{aligned}
\]
The right-hand side is the linear interpolation between the two adjacent
score-ordered prefix values.

On each interpolation interval the log divisor ratio is affine in
\(\log n\), while \(-\log\log\log n\) is strictly convex. One of the two
prefix endpoints therefore has at least as large a value of \(\log G\)
as any integer selection in the interval. Thus the supremum over integer
selections equals the supremum over score-ordered prefixes.

The total possible gain above the primorial is
\[
    -\sum_{p\leq p_m}\log(1-p^{-2})
    \leq
    \log\zeta(2).
\]
It follows that the prefix objective tends to minus infinity, so a
maximizing prefix exists. Comparing any maximizing prefix with the
primorial gives
\[
    \log\log\log N_m^\star
    \leq
    \log\log\theta(p_m)+\log\zeta(2),
\]
which is \eqref{e-prefix-bound}.

It remains to characterize an arbitrary maximizer \(N_m^\star\). The function
\(u\mapsto\log\log u\) is strictly concave, and its derivative at \(L_m\)
is \(\eps_m\). Optimality and the tangent inequality give, for every
integer \(M\) on the prescribed support,
\[
\begin{aligned}
    \log\frac{\sigma(M)}M-\log\log\log M
    &\leq
    \log\frac{\sigma(N_m^\star)}{N_m^\star}-\log\log L_m,\\
    \log\log\log M
    &\leq
    \log\log L_m+\eps_m(\log M-L_m).
\end{aligned}
\]
Consequently,
\[
    \log\frac{\sigma(M)}M-\eps_m\log M
    \leq
    \log\frac{\sigma(N_m^\star)}{N_m^\star}-\eps_m L_m.
\]
This expression is separable over the optional increments. Equation
\eqref{e-score-dictionary} shows that an increment is selected when its
score exceeds \(\eps_m\) and is unselected when its score is smaller.
Equality cannot occur. Adding or removing an increment with score
\(\eps_m\) would preserve the last display, while strictness in the
tangent inequality would produce a larger value of \(\log G\). Thus the
selected increments form a score-ordered prefix and satisfy
\eqref{e-integer-threshold}. This also proves the tie assertion and
extends \eqref{e-prefix-bound} to every maximizer.
\end{proof}

With exact score comparisons, the theorem gives a finite algorithm.
Maintain one score-ordered chain for each prime, repeatedly take
the largest available block, and evaluate \(G\) at the resulting prefix.
The search stops once \(\log N_j\) exceeds the bound in
\eqref{e-prefix-bound}. This computes the optimizer and its strict
threshold without searching an exponent box.

The threshold also gives the precise SA and CA interpretation. The
integer \(N_m^\star\) maximizes
\[
    \frac{\sigma(n)}{n^{1+\eps_m}}
\]
over integers on the prescribed support. It consequently maximizes
\(\sigma(n)/n\) on the same support under the size constraint
\(\log n\leq L_m\). It need not be a global CA or SA number because
the first-exponent increments \(F(p,1)\) are fixed by the support
constraint.

\subsection{The common scale}
\label{s-common-scale}

\begin{proposition}[Relaxed scale]
\label{prop-relaxed-scale}
As \(m\) tends to infinity,
\[
    S^\star\sim\theta(p_m)\sim p_m,
    \qquad
    T^\star\sim p_m\log p_m,
\]
and
\begin{equation}
\label{e-relaxed-displacement}
    S^\star-\theta(p_m)
    =
    O\left(\sqrt{\frac{p_m}{\log p_m}}\right).
\end{equation}
\end{proposition}

\begin{proof}
Put \(y=\sqrt{1+T^\star}\). We first show that \(y\leq p_m\) for all
sufficiently large \(m\). Otherwise every coordinate is interior, so
\[
    S^\star
    =
    2m\log y-\theta(p_m)
    \leq2m\log y.
\]
Since \(y^2-1=S^\star\log S^\star<2S^\star\log y\), this would imply
\(y/\log y\ll\sqrt m\). But \(y>p_m\) and the prime number theorem give
\(y/\log y>p_m/\log p_m\sim m\), a contradiction.

Splitting the response \eqref{e-relaxed-response} at \(p=y\) gives
\[
\begin{aligned}
    S^\star-\theta(p_m)
    &=
    2\bigl(\pi(y)\log y-\theta(y)\bigr)\\
    &=
    2\int_2^y\frac{\pi(u)}u\,du
    =
    O\left(\frac{y}{\log y}\right).
\end{aligned}
\]
Because \(y^2-1=S^\star\log S^\star\), the last bound is \(o(S^\star)\).
It follows that \(S^\star\sim\theta(p_m)\sim p_m\), and hence
\(T^\star\sim p_m\log p_m\). Substitution in the same bound proves
\eqref{e-relaxed-displacement}.
\end{proof}

\begin{proposition}[Integer scale]
\label{prop-integer-threshold-scale}
For every integer maximizer \(N_m^\star\),
\begin{equation}
\label{e-integer-mass}
    \log N_m^\star-\theta(p_m)
    =
    \bigl(\sqrt2+o(1)\bigr)
    \sqrt{\frac{T^\star}{\log T^\star}},
\end{equation}
and
\begin{equation}
\label{e-threshold-separation}
    1-T^\star\eps_m
    =
    \bigl(\sqrt2+o(1)\bigr)
    \sqrt{\frac{\log T^\star}{T^\star}}.
\end{equation}
In particular, \(T^\star\eps_m<1\) for all sufficiently large \(m\).
\end{proposition}

\begin{proof}
Put \(L_m=\log N_m^\star\). Strict separation in
Theorem~\ref{thm-integer-prefix} gives the exact identity
\[
    L_m-\theta(p_m)
    =
    \sum_{p\leq p_m}\log p\,
    \#\{k\geq2\mid F(p,k)>\eps_m\}.
\]
The bound \(F(p,k)\leq1/(p^k-1)\) first gives
\[
    L_m-\theta(p_m)
    =
    O\left(\sqrt{L_m\log L_m}\right).
\]
Since \(L_m=\theta(p_m)+O(\sqrt{L_m\log L_m})\), we have
\(L_m\sim\theta(p_m)\sim p_m\). Hence
\[
    L_m\log L_m\sim T^\star.
\]
We now estimate the score mass at this level. Put
\[
    U=L_m\log L_m.
\]
The first optional increment satisfies
\[
    F(p,2)
    =
    \frac{1+o(1)}{p^2\log p}.
\]
For every fixed \(\delta>0\) and every sufficiently large \(p\), this
gives the uniform bounds
\[
    \frac{1-\delta}{p^2\log p}
    \leq
    F(p,2)
    \leq
    \frac{1+\delta}{p^2\log p}.
\]
The positive solution of \(x^2\log x=U\) is
\[
    \bigl(\sqrt2+o(1)\bigr)
    \sqrt{\frac{U}{\log U}}.
\]
Comparison with the positive solutions of
\(x^2\log x=(1\pm\delta)U\), followed by the prime number theorem and
then \(\delta\to0\), gives
\[
    \sum_{\substack{p\leq p_m\\ F(p,2)>1/U}}\log p
    =
    \bigl(\sqrt2+o(1)\bigr)
    \sqrt{\frac{U}{\log U}}.
\]
The support cutoff is inactive because
\[
    p_m\asymp\frac{T^\star}{\log T^\star}
    \gg\sqrt{\frac{U}{\log U}}.
\]
For \(k\geq3\), the bound \(F(p,k)\leq1/(p^k-1)\) gives total log weight
\[
    \sum_{3\leq k\leq\log_2(1+U)}
    \theta\left((1+U)^{1/k}\right)
    =
    O(U^{1/3}+U^{1/4}\log U),
\]
which is lower order. Since \(U\sim T^\star\), the exact score identity
now proves \eqref{e-integer-mass}.

By \eqref{e-relaxed-displacement},
\[
    S^\star-\theta(p_m)
    =
    o\left(
        \sqrt{\frac{T^\star}{\log T^\star}}
    \right).
\]
Together with \eqref{e-integer-mass}, this gives
\[
    L_m-S^\star
    =
    \bigl(\sqrt2+o(1)\bigr)
    \sqrt{\frac{T^\star}{\log T^\star}}.
\]
Taylor expansion of \(u\log u\) now gives
\[
\begin{aligned}
    L_m\log L_m-T^\star
    &=
    (1+\log S^\star)(L_m-S^\star)
    +O\left(\frac{(L_m-S^\star)^2}{S^\star}\right)\\
    &=
    \bigl(\sqrt2+o(1)\bigr)
    \sqrt{T^\star\log T^\star}.
\end{aligned}
\]
Taking reciprocals proves \eqref{e-threshold-separation}.
\end{proof}

\section{Score rounding and the integrality gap}
\label{s-integrality-gap}

Fix an integer maximizer
\[
    N_m^\star=\prod_{i=1}^{m}p_i^{a_i},
\]
and put \(Y=\log(1+T^\star)\). We first show that its exponents are
obtained by rounding the relaxed exponents one coordinate at a time.

\begin{proposition}[One-step proximity]
\label{prop-floor-ceiling}
For every sufficiently large \(m\),
\[
    a_i
    \in
    \left\{
        \left\lfloor\frac{Y}{\log p_i}-1\right\rfloor,\,
        \left\lceil\frac{Y}{\log p_i}-1\right\rceil
    \right\}
\]
when \(p_i<e^{Y/2}\). For \(p_i\geq e^{Y/2}\), we have \(a_i=1\).
\end{proposition}

\begin{proof}
Fix a prime below \(e^{Y/2}\), and put
\[
    \ell=\left\lfloor\frac{Y}{\log p}-1\right\rfloor.
\]
The relaxed first-order equation is
\[
    \frac1{p^{Y/\log p}-1}=\frac1{T^\star}.
\]
If \(\ell\geq2\), the increment to exponent \(\ell\) lies below the
relaxed coordinate. Its score is greater than \(1/T^\star\), and hence greater than
\(\eps_m\) by Proposition~\ref{prop-integer-threshold-scale}. It is
selected. When \(\ell=1\), the support constraint already gives
\(a_i\geq\ell\).

We next rule out the increment to exponent \(\ell+2\). Write
\[
    u=\frac{Y}{\log p}-1+v
\]
on the corresponding unit interval. Since \(p^{Y/\log p}=1+T^\star\),
\[
    \frac{T^\star}{p^{u+1}-1}
    \leq p^{-v}.
\]
The interval begins at some \(v\in(0,1]\), so
\[
    T^\star F(p,\ell+2)
    \leq
    \int_0^1p^{-v}\,dv
    =
    \frac{1-p^{-1}}{\log p}
    \leq
    \frac1{2\log2}
    <1.
\]
Since \(T^\star\eps_m\) tends to one, this increment and all later ones
are unselected. Thus \(\ell\leq a_i\leq\ell+1\).

If \(Y/\log p-1\) is an integer, the same estimate applied to the next
increment shows that it is unselected. This proves the stated
floor--ceiling alternative. If \(p_i\geq e^{Y/2}\), apply the same
estimate to \(F(p_i,2)\). The corresponding unit interval begins at a
nonnegative value of \(v\), so this increment is unselected and
\(a_i=1\).
\end{proof}

\begin{theorem}[Integrality-gap asymptotic]
\label{thm-integrality-gap}
The integrality gap satisfies
\[
    \Igap(m)
    =
    \frac{2\sqrt2+o(1)}
         {\sqrt{p_m}\log p_m}.
\]
\end{theorem}

To compare the optimizers coordinatewise, define a score-rounded integer
exponent at every log-prime \(x\). For
\(Y>2\log2\) and \(\log2\leq x\leq Y/2\), put
\begin{equation}
\label{e-rounded-exponent}
    a_Y(x)
    =
    \begin{cases}
    \lfloor Y/x\rfloor,
        &F(e^x,\lfloor Y/x\rfloor)>(e^Y-1)^{-1},\\
    \lfloor Y/x\rfloor-1,
        &F(e^x,\lfloor Y/x\rfloor)\leq(e^Y-1)^{-1}.
    \end{cases}
\end{equation}
This is one of the two integers adjacent to \(Y/x-1\). It rounds upward
exactly when the intervening chord score exceeds the relaxed multiplier.
Define its local loss by
\begin{equation}
\label{e-kappa}
    \kappa(x,Y)
    =
    \log\frac{1-e^{-Y}}
              {1-e^{-x(a_Y(x)+1)}}
    +
    \frac{x(a_Y(x)+1)-Y}{e^Y-1},
\end{equation}
and define
\begin{equation}
\label{e-rounding-integral}
    J(Y)
    =
    \int_{\log2}^{Y/2}
    \kappa(x,Y)\frac{e^x}{x}\,dx.
\end{equation}
Thus \(\kappa\) is the loss in the local divisor factor after the
first-order log-size change has been removed.

\begin{proof}[Proof of Theorem~\ref{thm-integrality-gap}]
\medskip\noindent\emph{Rounding reduction.}
For \(p_i<e^{Y/2}\), let \(b_i=a_Y(\log p_i)\). Subtracting the
two objective values and linearizing the size penalty at \(S^\star\)
gives
\[
\begin{aligned}
    \Igap(m)
    &=
    \sum_{p_i<e^{Y/2}}
    \left[
        \log\frac{p_i-p_i^{-r_i^\star}}
                  {p_i-p_i^{-a_i}}
        -\frac{(r_i^\star-a_i)\log p_i}{T^\star}
    \right]\\
    &\quad
    +O\left(
        \frac{\log T^\star}{(T^\star)^2}
        (L_m-S^\star)^2
    \right).
\end{aligned}
\]
Proposition~\ref{prop-floor-ceiling} gives
\[
    |L_m-S^\star|
    \leq
    \theta(e^{Y/2})
    =
    O(\sqrt{T^\star}),
\]
so the displayed error is lower order on the required scale.

Both \(a_i\) and \(b_i\) are adjacent to \(r_i^\star\). If they differ,
the intervening score lies between \(\eps_m\) and \(1/T^\star\). Hence
replacing \(a_i\) by \(b_i\) in the sum changes it by at most
\[
    \theta(e^{Y/2})
    \left(\frac1{T^\star}-\eps_m\right)
    =
    O\left(\frac{\sqrt{\log T^\star}}{T^\star}\right),
\]
where we used \eqref{e-threshold-separation}. This is also lower order.
The bracket with \(b_i\) is exactly
\(\kappa(\log p_i,Y)\), by \eqref{e-kappa}.
We have proved
\[
    \sum_{p<e^{Y/2}}\kappa(\log p,Y)
    =
    \Igap(m)
    +
    o\left(\frac1{\sqrt{T^\star\log T^\star}}\right).
\]

\medskip\noindent\emph{Prime-sum approximation.}
Put \(q=e^{-Y}\) and \(\lambda=(e^Y-1)^{-1}\). On an interval where
\(k=\lfloor Y/x\rfloor\) is constant, put \(s=Y-kx\). Formula
\eqref{e-kappa} becomes
\begin{equation}
\label{e-kappa-branches}
    \kappa(x,Y)
    =
    \begin{cases}
    \log(1-q)-\log(1-qe^s)-s\lambda,
        &F(e^x,k)\leq\lambda,\\[1.5mm]
    \log(1-q)-\log(1-qe^{s-x})+(x-s)\lambda,
        &F(e^x,k)>\lambda.
    \end{cases}
\end{equation}
The difference between the branches is
\[
    x\bigl(\lambda-F(e^x,k)\bigr),
\]
so they agree at a switch.

On the first branch,
\[
\begin{aligned}
    xF(e^x,k)
    &=
    \log(1-qe^{s-x})-\log(1-qe^s)\\
    &\geq
    qe^s(1-e^{-x})
    \geq
    \tfrac12qe^s.
\end{aligned}
\]
The branch condition therefore gives \(qe^s\ll xq\). On the second
branch, \(qe^{s-x}\leq q\). It follows from
\eqref{e-kappa-branches} that
\[
    |\kappa(x,Y)|\ll Y e^{-Y}.
\]
Within a shell, \(s'=-k\), and differentiation gives
\[
    |\partial_x\kappa(x,Y)|\ll Y e^{-Y}
\]
on every smooth piece. There are \(O(Y)\) shells. The score
\[
    F(e^x,k)=\int_k^{k+1}\frac{du}{e^{xu}-1}
\]
is \eqref{e-score-integral} after shifting the integration variable. It
decreases strictly with \(x\), so each shell has at most one branch
switch. The smooth pieces have total length \(O(Y)\), and the
shell-boundary jumps are \(O(Ye^{-Y})\). Consequently, the total
variation satisfies
\[
    \operatorname{TV}_{[\log2,Y/2]}
    \kappa(\,\cdot\,,Y)
    \ll
    Y^2e^{-Y}.
\]
At \(x=Y/2\), the score-rounded exponent is one and
\(\kappa(Y/2,Y)=0\).

The de la Vall\'ee Poussin estimate \cite{delavalleepoussin1899} gives
\[
    \pi(e^x)-\operatorname{li}(e^x)
    =
    O\left(
        \frac{e^x}{x}e^{-c\sqrt x}
    \right).
\]
Stieltjes integration by parts, using the variation bound and the
vanishing upper endpoint, now gives
\[
\begin{aligned}
    &\sum_{p<e^{Y/2}}\kappa(\log p,Y)
    -
    \int_{\log2}^{Y/2}\kappa(x,Y)\frac{e^x}{x}\,dx\\
    &\qquad
    =
    O\left(
        Ye^{-Y/2}e^{-c'\sqrt Y}
    \right)
    +O(Ye^{-Y}).
\end{aligned}
\]
This is lower order than the required scale. Shell-boundary conventions
change the expression by at most \(O(Y^2e^{-Y})\).
It follows that
\[
    \sum_{p<e^{Y/2}}\kappa(\log p,Y)
    =
    J(Y)
    +
    o\left(\frac1{\sqrt{(e^Y-1)Y}}\right).
\]

\medskip\noindent\emph{Second-exponent range.}
The branch formulas give the sharper bound
\[
    |\kappa(x,Y)|\ll q(1+x).
\]
Therefore
\[
    \int_{\log2}^{Y/3}
    |\kappa(x,Y)|\frac{e^x}{x}\,dx
    =
    O(e^{-2Y/3}),
\]
which is lower order. Only \(Y/3<x\leq Y/2\), where
\(\lfloor Y/x\rfloor=2\), contributes at first order.

For all sufficiently large \(Y\), let \(x_\star\) be the solution in
\((Y/3,Y/2)\) of
\begin{equation}
\label{e-rounding-switch}
    F(e^{x_\star},2)=\frac1{e^Y-1}.
\end{equation}
On this shell put
\[
    u=Y-2x,
    \qquad
    x=\frac{Y-u}{2}.
\]
The score in \eqref{e-rounding-switch} increases with \(u\). It is below
\((e^Y-1)^{-1}\) at \(u=0\) and above it near \(u=Y/3\), so the switch
exists and is unique. Put \(u_\star=Y-2x_\star\). Uniformly on the shell,
\[
\begin{aligned}
    xF(e^x,2)
    &=
    \log(1-e^{-3x})-\log(1-e^{-2x})\\
    &=
    e^{-Y+u}
    \bigl(1+O(e^{-x}+e^{-Y+u})\bigr).
\end{aligned}
\]
The switch equation gives
\[
    e^{u_\star}
    =
    x_\star\bigl(1+O(e^{-Y/3})\bigr).
\]
Together with \(2x_\star+u_\star=Y\), this proves
\[
    2x_\star+\log x_\star=Y+o(1).
\]

On the two sides of the switch, \eqref{e-kappa-branches} gives
\[
    \kappa(x,Y)
    =
    \begin{cases}
    q(e^u-1-u)+O(q^2e^{2u}),
        &0\leq u\leq u_\star,\\[1mm]
    q(x-u-1+e^{-(x-u)})+O(Yq^2),
        &u_\star\leq u<Y/3.
\end{cases}
\]
Since \(dx=-du/2\) and \(e^x=e^{Y/2}e^{-u/2}\), the contribution from
each branch is \(qe^{Y/2}/2\) times its corresponding integral below.
For the first branch,
\[
    \int_0^{u_\star}
    \frac{e^{-u/2}}x(e^u-1-u)\,du
    =
    \frac{2+o(1)}{\sqrt{x_\star}}.
\]
Indeed, the terms \(1+u\) contribute \(O(x_\star^{-1})\). In the
remaining term, set \(t=u_\star-u\), use
\(x=x_\star+t/2\) and \(e^{u_\star}\sim x_\star\), and apply dominated
convergence.

For the second branch, set \(t=u-u_\star\). Then
\[
    x=x_\star-\frac t2,
    \qquad
    x-u=x_\star-u_\star-\frac{3t}{2}.
\]
Since
\[
    0\leq x-u-1+e^{-(x-u)}\leq x-u\leq x,
\]
the second integral equals
\[
    e^{-u_\star/2}
    \int_0^{Y/3-u_\star}
    e^{-t/2}
    \frac{x-u-1+e^{-(x-u)}}x\,dt.
\]
The ratio is bounded by one and tends to one for every fixed \(t\).
Since \(e^{-u_\star/2}\sim x_\star^{-1/2}\), dominated convergence gives
\[
    \int_{u_\star}^{Y/3}
    \frac{e^{-u/2}}x
    \bigl(x-u-1+e^{-(x-u)}\bigr)\,du
    =
    \frac{2+o(1)}{\sqrt{x_\star}}.
\]
After multiplication by the Jacobian factor \(e^{Y/2}/2\), the error
\(O(q^2e^{2u})\) on the first branch contributes
\[
    O(e^{-3Y/2}x_\star^{1/2}),
\]
while the error \(O(Yq^2)\) on the second branch contributes
\[
    O(e^{-3Y/2}x_\star^{-1/2}).
\]
Both are \(o(e^{-Y/2}x_\star^{-1/2})\), the scale of the main term.
Combining the two branches yields
\[
    J(Y)
    =
    e^{-Y/2}
    \left(
        \frac2{\sqrt{x_\star}}
        +o(x_\star^{-1/2})
    \right).
\]
Finally, \(x_\star\sim Y/2\) and
\[
    e^{-Y/2}
    =
    \frac{\sqrt{1-e^{-Y}}}{\sqrt{e^Y-1}}.
\]
Consequently,
\[
    J(Y)
    =
    \frac1{\sqrt{e^Y-1}}
    \left(
        \frac2{\sqrt{x_\star}}
        +o(x_\star^{-1/2})
    \right)
    =
    \frac{2\sqrt2+o(1)}
         {\sqrt{(e^Y-1)Y}}.
\]
The two sides of the switch contribute equally at first order. Since
\(e^Y-1=T^\star\), the rounding reduction and prime-sum approximation
give
\[
    \Igap(m)
    =
    \frac{2\sqrt2+o(1)}
         {\sqrt{T^\star Y}}.
\]
Since \(Y\sim\log T^\star\) and
\(T^\star\sim p_m\log p_m\), the denominator is asymptotic to
\(\sqrt{p_m}\log p_m\).
\end{proof}

\subsection{Relation to CA shells and the Axler--Nicolas comparison}
\label{s-score-connection}

The fixed-support optimizer is a constrained CA point, since its
optional increments are selected by the same score \(F\). It is also a
size-constrained record on its prescribed support. The classical CA shells
therefore give a second description of the integer optimizer, although
they do not provide the relaxed response or the rounding loss used in the
proof above.

Define the integer shell thresholds by
\[
    F(\xi,1)=\eps_m,
    \qquad
    F(\xi_k,k)=\eps_m.
\]
Nicolas \cite[Proposition~4.7]{nicolas-2022-sigma} proves
\[
    \xi^{1/k}\leq\xi_k\leq(k\xi)^{1/k},
    \qquad
    \xi_2\sim\sqrt{2\xi}.
\]
Since
\[
    F(\xi,1)\sim\frac1{\xi\log\xi},
\]
Proposition~\ref{prop-integer-threshold-scale} gives
\(\xi\sim L_m\sim p_m\). The exact threshold rule is
\[
    L_m-\theta(p_m)
    =
    \sum_{k\geq2}\theta\bigl(\min\{p_m,\xi_k\}\bigr).
\]
Nicolas's bounds make every support cap on the right inactive for all
sufficiently large \(m\). They then give
\[
    L_m-\theta(p_m)
    =
    \sum_{k\geq2}\theta(\xi_k)
    =
    \theta(\xi_2)+O(\xi^{1/3})
    =
    \sqrt{2p_m}+o(\sqrt{p_m}).
\]
The middle estimate follows from Nicolas's bounds for \(\xi_k\), the
prime number theorem, and the fact that only \(O(\log\xi)\) shells are
nonempty.

The relaxed switch is nearby but not identical. It is defined by
\[
    F(e^{x_\star},2)=\frac1{T^\star},
\]
whereas \(\xi_2\) uses \(\eps_m\). Equation
\eqref{e-threshold-separation} shows that
\[
    e^{x_\star}\sim\xi_2\sim\sqrt{2p_m}.
\]
Thus the standard CA second shell locates the score-rounding
transition, while the prime-sum integral resolves the loss on its two
sides.

The exact integer factorization, relative to the Mertens term used in
the next section, has the correction
\[
    \sum_{i=1}^m\log(1-p_i^{-(a_i+1)})
    -\bigl(\log\log L_m-\log\log\theta(p_m)\bigr).
\]
For \(\xi_2<p\leq p_m\), the exponent of \(p\) is one. The prime-square
tail estimate \cite[equations~(2.10)--(2.11)]{nicolas-2022-sigma} gives
\[
    \sum_{\xi_2<p\leq p_m}\log(1-p^{-2})
    =
    -\frac1{\xi_2\log\xi_2}
    +o\left(\frac1{\sqrt{p_m}\log p_m}\right)
    =
    -\frac{\sqrt2+o(1)}{\sqrt{p_m}\log p_m}.
\]
The remaining Euler factors are lower order. Primes with exponent two
lie above \(\xi_3\), so their contribution is
\[
    O\left(\sum_{p>\xi_3}p^{-3}\right)
    =
    O(\xi^{-2/3}).
\]
If \(a_i\geq3\), the next increment is unselected. Hence
\(p_i>\xi_{a_i+1}\), and Nicolas's lower bound gives
\(p_i^{a_i+1}>\xi\). There are \(O(\xi_3)=O(\xi^{1/3})\) such primes, so
their total contribution is also \(O(\xi^{-2/3})\). Consequently,
\[
    \sum_{i=1}^m\log(1-p_i^{-(a_i+1)})
    =
    -\frac{\sqrt2+o(1)}{\sqrt{p_m}\log p_m}.
\]
Taylor expansion and the shell-size formula similarly give
\[
    \log\log L_m-\log\log\theta(p_m)
    =
    \frac{\sqrt2+o(1)}{\sqrt{p_m}\log p_m}.
\]
The CA decomposition partitions the correction into an Euler-factor tail
and a log-size change, whereas the rounding integral partitions it at the
downward and upward rounding switch.

Nicolas's expansion of the Axler--Nicolas size-constrained comparison
\cite[equations~(5.3)--(5.5)]{nicolas-2025} contains the same formal pair
\[
    \frac1{\xi_2\log\xi_2},
    \qquad
    \frac{\xi_2}{\xi\log\xi}.
\]
The first is the same leading Euler-factor tail. Their second term measures
additional support-edge primes under a common size constraint. Ours measures
the Robin penalty applied to second-shell log-size. The common score and
shell equation explain the same \(2\sqrt2\) coefficient. The different
feasible sets and the different second terms prevent either theorem from
being a corollary of the other.

Since the Axler--Nicolas ratio tends to one, its relative gap and its
logarithm have the same leading term. The normalizations also agree after
the natural endpoint change. If \(P\)
is the largest prime in the primorial below \(X\), then
\(\log X\sim\theta(P)\sim P\). Hence the Axler--Nicolas scale satisfies
\[
    \frac1{\sqrt{\log X}\log\log X}
    \sim
    \frac1{\sqrt P\log P},
\]
which is the fixed-slice scale when \(P=p_m\).

There is a related relaxation interpretation. For every prime \(p\),
\[
    \frac{\sigma(p^a)}{p^a}
    \longrightarrow
    \frac1{1-p^{-1}}
    =
    \frac{p^a}{\varphi(p^a)}
\]
as \(a\) tends to infinity. Their ratio can therefore be viewed as
comparing the infinite-exponent local envelope with finite integer
exponents while also optimizing the support under a size constraint. The
present problem compares finite real and integer exponents on a fixed
support. In both
cases the first nontrivial marginal decision is the transition from
exponent one to exponent two. This common decision, rather than a formal
reduction between the problems, accounts for the matching constant.

\section{The analytic term}
\label{s-relaxed-excess}

Define the Chebyshev-normalized Mertens error
\begin{equation}
\label{e-mertens-error}
    E_\theta(x)
    =
    -\sum_{p\leq x}\log\left(1-\frac1p\right)
    -\gamma-\log\log\theta(x).
\end{equation}
The relaxed problem differs from this classical term by less than the
integrality gap.

\paragraph{Relaxed value.}
We first show that
\begin{equation}
\label{e-relaxed-value}
    g_{\reals}(m)
    =
    E_\theta(p_m)
    +
    O\left(
        \frac1{\sqrt{p_m}\log^{3/2}p_m}
    \right).
\end{equation}
Put \(y=\sqrt{1+T^\star}\). The response
\eqref{e-relaxed-response} gives
\[
    p_i-p_i^{-r_i^\star}
    =
    \frac{p_iT^\star}{1+T^\star}
\]
for \(p_i\leq y\), while \(r_i^\star=1\) for \(p_i>y\).
Substitution in \eqref{e-objective} gives
\[
\begin{aligned}
    g_{\reals}(m)
    &=
    \pi(y)\log\frac{T^\star}{1+T^\star}
    -\sum_{p\leq p_m}\log\left(1-\frac1p\right)\\
    &\quad
    +\sum_{y<p\leq p_m}\log\left(1-\frac1{p^2}\right)
    -\gamma-\log\log S^\star.
\end{aligned}
\]
The prime number theorem and partial summation give
\[
    \pi(y)\left|\log\frac{T^\star}{1+T^\star}\right|
    +
    \sum_{p>y}\left|\log\left(1-\frac1{p^2}\right)\right|
    =
    O\left(\frac1{\sqrt{T^\star}\log T^\star}\right).
\]
Proposition~\ref{prop-relaxed-scale} and
\eqref{e-relaxed-displacement} also give
\[
    \log\log S^\star-\log\log\theta(p_m)
    =
    O\left(\frac1{\sqrt{T^\star}\log T^\star}\right).
\]
The result follows from \(T^\star\sim p_m\log p_m\).

We now use Nicolas's results to evaluate \(E_\theta(p_m)\).

\subsection{Under the Riemann hypothesis}
\label{s-rh-branch}

Assume the Riemann hypothesis, and define
\[
    \mathcal H(t)
    =
    \sum_\rho
    \frac{e^{(\rho-1/2)t}}{\rho(\rho-1)},
    \qquad
    \tau=2+\gamma-\log(4\pi),
\]
where the sum runs over the nontrivial zeros of the zeta function, with
multiplicity. The series is absolutely convergent because its terms are
\(O(|\operatorname{Im}\rho|^{-2})\).

Nicolas's integrated explicit formula
\cite[Theorem~3(b), final formula in its proof]
{nicolas-1983-petites-valeurs}, translated to the normalization
\eqref{e-mertens-error}, is
\begin{equation}
\label{e-nicolas-integrated}
    E_\theta(x)
    =
    \frac{2-\mathcal H(\log x)}
         {\sqrt x\log x}
    +
    O\left(\frac{1}{\sqrt x\log^2x}\right).
\end{equation}
The same proof records the classical identity
\[
    \sum_\rho\frac{1}{\rho(1-\rho)}
    =
    2+\gamma-\log(4\pi).
\]
On the critical line, \(\rho(1-\rho)=|\rho|^2\). Therefore
\begin{equation}
\label{e-H-bound}
    |\mathcal H(t)|\leq\tau.
\end{equation}

\begin{theorem}[Relaxed excess under the Riemann hypothesis]
\label{thm-rh-relaxed}
Assume the Riemann hypothesis. Then
\[
    \sqrt{p_m}\log p_m\,g_{\reals}(m)
    =
    2-\mathcal H(\log p_m)
    +O\left(\frac1{\sqrt{\log p_m}}\right).
\]
In particular,
\[
    2-\tau+o(1)
    \leq
    \sqrt{p_m}\log p_m\,g_{\reals}(m)
    \leq
    2+\tau+o(1),
\]
and \(g_{\reals}(m)>0\) for all sufficiently large \(m\).
\end{theorem}

\begin{proof}
Apply \eqref{e-nicolas-integrated} at \(x=p_m\), and use
\eqref{e-relaxed-value}. The error in \eqref{e-relaxed-value}, after
multiplication by \(\sqrt{p_m}\log p_m\), is
\(O(1/\sqrt{\log p_m})\). The bounds follow from
\eqref{e-H-bound}. The lower one is eventually positive because
\(2-\tau>0\).
\end{proof}

The bound describes a band containing the normalized relaxed excess. It
does not assert that the full band is its limit set or that either
endpoint is attained along the primes.

\subsection{Off the Riemann hypothesis}
\label{s-off-rh-branch}

We use \(f=\Omega_\pm(g)\) when \(f/g\) has a positive limsup and a
negative liminf.

\begin{theorem}[Off-RH excursions]
\label{thm-off-rh-relaxed}
Suppose the Riemann hypothesis is false, and let
\[
    \Theta
    =
    \sup\{\operatorname{Re}\rho\mid\zeta(\rho)=0\}.
\]
For every \(0<\eps<\Theta-1/2\),
\[
    g_{\reals}(m)
    =
    \Omega_\pm\left(p_m^{-1/2+\eps}\right).
\]
Consequently,
\[
\begin{aligned}
    \limsup_{m\to\infty}
    \sqrt{p_m}\log p_m\,g_{\reals}(m)
    &=+\infty,\\
    \liminf_{m\to\infty}
    \sqrt{p_m}\log p_m\,g_{\reals}(m)
    &=-\infty.
\end{aligned}
\]
\end{theorem}

\begin{proof}
Nicolas
\cite[Theorem~3(c)]{nicolas-1983-petites-valeurs} proves, in the
normalization \eqref{e-mertens-error}, that
\[
    E_\theta(x)=\Omega_\pm(x^{-b})
    \qquad
    (1-\Theta<b<1/2).
\]
The function \(E_\theta\) is constant between consecutive primes, and
the prime number theorem makes the preceding prime asymptotic to \(x\).
The two real sequences in Nicolas's result therefore give two prime
subsequences with the same orders and opposite signs.

By \eqref{e-relaxed-value},
\[
    g_{\reals}(m)
    =
    E_\theta(p_m)
    +
    O\left(
        \frac{1}{\sqrt{p_m}\log^{3/2}p_m}
    \right).
\]
The error is \(o(p_m^{-b})\). Taking \(b=1/2-\eps\) proves the first
claim, and the normalized divergence follows.
\end{proof}

\section{Consequences for Robin's inequality}
\label{s-robin-margin}

The two preceding results, together with
\(-g_{\integers}(m)=\Igap(m)-g_{\reals}(m)\), give the integer margin.

\begin{theorem}[Fixed-slice Robin margin]
\label{thm-rh-margin}
Under the Riemann hypothesis,
\begin{equation}
\label{e-rh-margin}
    -\sqrt{p_m}\log p_m\,g_{\integers}(m)
    =
    2\sqrt2-2+\mathcal H(\log p_m)+o(1).
\end{equation}
Consequently,
\[
    g_{\reals}(m)>0>g_{\integers}(m)
\]
for every sufficiently large \(m\).

If the Riemann hypothesis is false, then
\[
\begin{aligned}
    \limsup_{m\to\infty}
    -\sqrt{p_m}\log p_m\,g_{\integers}(m)
    &=+\infty,\\
    \liminf_{m\to\infty}
    -\sqrt{p_m}\log p_m\,g_{\integers}(m)
    &=-\infty.
\end{aligned}
\]
\end{theorem}

\begin{proof}
Under the Riemann hypothesis, subtract
Theorem~\ref{thm-rh-relaxed} from the gap formula in
Theorem~\ref{thm-integrality-gap}. This gives \eqref{e-rh-margin}.
The relaxed value is eventually positive by
Theorem~\ref{thm-rh-relaxed}. The integer value is eventually negative
because
\[
    2\sqrt2-2+\mathcal H(\log p_m)
    \geq
    2\sqrt2-2-\tau
    >
    0.
\]

If the hypothesis is false, Theorem~\ref{thm-off-rh-relaxed} gives
opposite-sign excursions of size \(p_m^{-b}\) for every admissible
\(b<1/2\). The integrality gap is
\[
    O\left(\frac1{\sqrt{p_m}\log p_m}\right)
    =
    o(p_m^{-b}).
\]
It cannot suppress either excursion. Since
\(-g_{\integers}=\Igap-g_{\reals}\), the margin has the claimed
opposite-sign divergence.
\end{proof}

\begin{corollary}[Eventual fixed-\texorpdfstring{\(\omega\)}{omega}
criterion]
\label{cor-fixed-omega-equivalence}
The Riemann hypothesis is equivalent to the existence of an \(M\) such
that Robin's inequality holds for every \(n\geq5041\) with
\(\omega(n)=m\), for every \(m\geq M\).
\end{corollary}

\begin{proof}
If the Riemann hypothesis holds, the theorem gives
\(g_{\integers}(m)<0\) for all sufficiently large \(m\), so every
integer in each such slice satisfies Robin's inequality.

If it fails, the theorem gives infinitely many \(m\) with
\(g_{\integers}(m)>0\). Proposition~\ref{prop-support} shows that the
maximum is attained, so each such slice contains a violation.
\end{proof}

Formula \eqref{e-rh-margin} separates the sign quantitatively. The
analytic term contributes \(2-\mathcal H\), while the rounding
correction contributes \(2\sqrt2\). The strict inequality
\(2\sqrt2>2+\tau\) is why Robin's inequality survives on the integer
slice but fails for its real-exponent relaxation.

\section{At most \texorpdfstring{\(m\)}{m} prime factors}
\label{s-bounded-support}

When \(\omega(n)\leq m\), the support is chosen together with the
exponents. Define
\[
    g_{\integers}^{\leq}(m)
    =
    \sup_{\substack{n\geq5041\\ \omega(n)\leq m}}
    \log G(n).
\]
The CA marginal ordering restricted to the first \(m\) primes solves this
problem, apart from a finite boundary created by \(n\geq5041\).

For \(p\in\{p_1,\ldots,p_m\}\), order all increments \(F(p,k)\),
\(k\geq1\), by decreasing score, processing equal scores as one block.
Starting from \(1\), let \(n_j\) be the integer obtained after the first
\(j\) blocks. Every prefix is a valid exponent vector with at
most \(m\) positive exponents. Direct score comparison gives the first
score-ordered prefixes
\[
    2,\ 6,\ 12,\ 60,\ 120,\ 360,\ 2520,\ 5040,\ 55440,
\]
which agree with Nicolas's CA tables through \(2520\) and at \(5040\) and
\(55440\) \cite[Tables~1--2]{nicolas-2022-sigma}. Thus \(55440\) is the
first score-ordered prefix above Robin's lower cutoff.

\begin{theorem}[Bounded-support characterization]
\label{thm-bounded-support-prefix}
For every \(m\geq6\),
\begin{equation}
\label{e-bounded-support-prefix}
    g_{\integers}^{\leq}(m)
    =
    \max\left\{
        \max_{\substack{5041\leq n<55440\\ \omega(n)\leq m}}
        \log G(n),
        \ 
        \max_{\substack{j\geq1\\ n_j\geq55440}}
        \log G(n_j)
    \right\}.
\end{equation}
Both maxima are attained. Every maximizing integer at least \(55440\)
is a score-ordered prefix.
\end{theorem}

\begin{proof}
We first reduce an arbitrary support to an initial prime segment. Suppose
\[
    n=\prod_{i=1}^{k}q_i^{a_i},
    \qquad
    q_1<\cdots<q_k,
    \qquad
    k\leq m.
\]
Replacing \(q_i\) by \(p_i\) decreases \(n\) and strictly increases
\(\sigma(n)/n\), unless the support was already initial. If the resulting
integer is below \(5041\), multiply it by the smallest power of \(2\)
that reaches \(5041\). The result is below \(10082\), preserves the
support size, and further increases the divisor ratio. For
\(n\geq55440\), this produces a better feasible integer, either in the
finite range in \eqref{e-bounded-support-prefix} or on an initial
support.

An integer on an initial support is a finite selection from the ordered
increment list. Put \(z=\log n\), and choose consecutive prefix integers
with
\[
    \log n_j\leq z\leq\log n_{j+1}.
\]
The fractional-knapsack inequality bounds
\(\log(\sigma(n)/n)\) by the chord joining the corresponding prefix
values. This chord is affine in \(z\), while \(-\log\log z\) is strictly
convex. One of its endpoints therefore has objective at least
\(\log G(n)\). If \(z=\log55440\), then \(n=55440\), which is already a
prefix. If \(z>\log55440\), both adjacent endpoints are feasible. Equality
for a maximizing integer then requires equality in the knapsack bound. If
the selection cuts a tie block, its point lies in the interior of the
affine block segment, where strict convexity makes one endpoint strictly
better.
At an endpoint, equality in the knapsack bound forces every earlier block
to be selected and every later block to be unselected. The total log-size
then forces the intervening block to be either complete or absent. Thus
every maximizing integer in this range is a score-ordered prefix.

The first maximum in \eqref{e-bounded-support-prefix} is finite. Along
the prefixes, the total possible increase in the log divisor ratio is
bounded by
\[
    -\sum_{i=1}^{m}\log\left(1-\frac1{p_i}\right),
\]
whereas \(\log\log\log n_j\) tends to infinity. Hence
\(\log G(n_j)\) tends to minus infinity, and the second maximum is also
attained.
\end{proof}

\begin{corollary}[Unsaturated support]
\label{cor-bounded-support-ca}
Let \(n>55440\) maximize \(g_{\integers}^{\leq}(m)\), and suppose
\(k=\omega(n)<m\). Then \(n\) is the unique colossally abundant number
with parameter
\[
    \eps=\frac1{\log n\log\log n}.
\]
It is also superabundant at its own size, and
\[
    p_k<\log n<p_{k+1}+1,
    \qquad
    k\in\{\pi(\log n)-1,\pi(\log n)\}.
\]
\end{corollary}

\begin{proof}
Let the last selected block have score \(c_-\), and let the next block
have score \(c_+\). On the two adjacent chord intervals, the one-sided
derivatives of the interpolated objective at \(\log n\) are
\[
    c_- -\eps>0,
    \qquad
    c_+ -\eps<0.
\]
The weak inequalities follow from endpoint maximality, and equality is
excluded by strict convexity on the adjacent interval. Thus every
selected score exceeds \(\eps\), while every eligible unselected score is
smaller.

Because \(k<m\), the first-exponent increment for \(p_{k+1}\) is eligible
and unselected.
Monotonicity of \(F\) extends the strict separation to every larger prime
and to all its exponent increments. Equation~\eqref{e-score-dictionary}
then shows that \(n\) uniquely maximizes
\[
    \frac{\sigma(N)}{N^{1+\eps}}
\]
over all positive integers, which is the CA characterization. The same
inequality shows that \(\sigma(N)/N\leq\sigma(n)/n\) whenever \(N\leq n\),
so \(n\) is SA.

The selected first-exponent increment at \(p_k\) and the unselected one
at \(p_{k+1}\) give
\[
    F(p_k,1)>\frac1{\log n\log\log n}>F(p_{k+1},1).
\]
Using
\[
    \frac1{p+1}
    <
    \log\left(1+\frac1p\right)
    <
    \frac1p
\]
and the monotonicity of \(x\log x\) gives the endpoint bounds. The
statement about \(k\) follows.
\end{proof}

The two maxima in \eqref{e-bounded-support-prefix} exchange roles as
\(m\) grows.

\begin{proposition}[Eventual prefix dominance]
\label{prop-eventual-prefix}
For every \(m\geq6\), the first maximum in
\eqref{e-bounded-support-prefix} equals
\(\log G(10080)=-0.0142\ldots\). For \(6\leq m\leq29\), it is the
larger one, \(g_{\integers}^{\leq}(m)=\log G(10080)\), and the unique
maximizer is \(10080=2^{5}3^{2}5\cdot7\). For \(m\geq30\), every
maximizer of \(g_{\integers}^{\leq}(m)\) is a score-ordered prefix
above \(55440\).
\end{proposition}

\begin{proof}
Every integer below \(55440\) has at most six distinct prime factors,
so the first maximum does not depend on \(m\geq6\). A verification over
the window \([5041,55440)\) places it at \(10080\), with the runner-up
\(27720\) lower by \(0.0075\ldots\).

A score-ordered prefix \(n_j\geq55440\) with \(k\) positive exponents
lies in the slice \(\omega(n)=k\), so its objective is at most
\(g_{\integers}(k)\), with equality only at a slice maximizer. For
\(k\leq5\), Proposition~\ref{prop-small-slices} gives
\(g_{\integers}(k)\leq g_{\integers}(4)=\log G(10080)\), with equality
only at \(k=4\), whose maximizer \(10080\) lies below \(55440\), so the
prefix objective is strictly smaller. For \(6\leq k\leq29\), a
verification by the prefix computation of
Theorem~\ref{thm-integer-prefix} gives
\(g_{\integers}(k)<\log G(10080)\). The tightest case is
\(g_{\integers}(29)=-0.0144\ldots\), short by \(0.00012\ldots\). Hence
the second maximum stays strictly below the first for
\(6\leq m\leq29\), which proves the middle claim.

The same computation gives
\(g_{\integers}(30)=-0.0135\ldots>\log G(10080)\). For \(m\geq30\) the
slice \(\omega(n)=30\) is feasible, so
\(g_{\integers}^{\leq}(m)>\log G(10080)\), no integer in the window
attains the maximum, and Theorem~\ref{thm-bounded-support-prefix}
identifies every maximizer as a score-ordered prefix.
\end{proof}

The prefixes must eventually dominate. The slice maxima
\(g_{\integers}(m)\) tend to zero, since the relaxed excess is a
Mertens error by \eqref{e-relaxed-value} and the gap vanishes by
Theorem~\ref{thm-integrality-gap}, while \(\log G(10080)\) is a fixed
negative number. The computation locates the change at \(m=30\).

For a maximizing prefix above \(55440\), the threshold has two
interpretations. If the cap is active, the Lagrangian and size-constrained
statements hold over integers whose prime divisors are at most \(p_m\). If
the cap is inactive, strict separation extends to all primes, so the
optimizer is an unrestricted CA number and hence SA. Thus passing from
\(\omega(n)=m\) to \(\omega(n)\leq m\) restores the first-exponent scores
\(F(p,1)\) as decisions.

\section{Conclusion}
\label{s-conclusion}

The fixed-\(\omega\) problem has complementary one-dimensional
descriptions. A scalar equation gives the relaxed optimizer, while
ordering increments by Nicolas's score gives every integer optimizer as a
score-ordered prefix and supplies a finite computation. The integer exponents
are within one step of the relaxed exponents. Transferring their rounding
loss to a prime integral gives
\(2\sqrt2/(\sqrt{p_m}\log p_m)\), with equal contributions from the two
sides of the second-exponent switch.

Classical CA shells locate the same switch and explain the common
coefficient in the Axler--Nicolas size-constrained comparison, although
the feasible sets and the second terms differ. Combining the unconditional
gap with Nicolas's Mertens formula gives the eventual fixed-\(\omega\)
criterion. The same score ordering also characterizes the bounded-support
problem \(\omega(n)\leq m\).


\bibliographystyle{plain}
\bibliography{refs}

\end{document}
